\chapter{Thắng trong mắt người đời}
\markboth{Thắng trong mắt người đời}{Winning in Public Eyes}

\section{Một học sinh thắng cuộc tranh luận nhưng mất người bạn thân.}
\begin{truyen}{Một học sinh thắng cuộc tranh luận nhưng mất người bạn thân.}{Win the argument, lose the friend.}
\chuhoa{E}{m nói hay, lý lẽ chặt, cả lớp im lặng vì không cãi lại được. Nhưng sau giờ học, người bạn kia không còn nói chuyện như trước nữa. Em thắng cuộc tranh luận, nhưng thua một mối thân tình.}
\textit{(The student spoke sharply and no one in class could answer back. But after class, the friend no longer spoke the same way. The argument was won, but a friendship was lost.)}
\end{truyen}
\begin{baihoc}
Không phải cái thắng nào cũng đáng vui nếu cái giá là một con người.
\textit{(Not every victory is worth celebrating if the price is a person.)}
\end{baihoc}
\begin{ghinhoanh}\textbf{Short English to remember:}

Win the argument, lose the friend.
\end{ghinhoanh}
\begin{tuhocnhanh}
1. Em có đang cố thắng ai bằng lời nói không? \textit{(Are you trying to defeat someone with words right now?)}

2. Điều gì đáng giữ hơn: phần đúng hay người đối diện? \textit{(What is worth keeping more: being right or keeping the person?)}
\end{tuhocnhanh}
\ngancach

\section{Một người khoe thành tích đến mức ai cũng mệt.}
\begin{truyen}{Một người khoe thành tích đến mức ai cũng mệt.}{Pride can empty applause.}
\chuhoa{A}{nh thành công thật, nhưng đi đâu cũng nhắc để được công nhận. Ban đầu người ta chúc mừng, dần dần người ta tránh. Anh có thêm tiếng tăm, nhưng mất sự dễ chịu trong lòng người khác.}
\textit{(He truly succeeded, but mentioned it everywhere to gain recognition. At first people congratulated him; later they avoided him. He gained status but lost ease in other people's hearts.)}
\end{truyen}
\begin{baihoc}
Khoe quá nhiều thường làm chiến thắng nhỏ đi trong mắt người khác.
\textit{(Showing off too much often makes a victory feel smaller in other people's eyes.)}
\end{baihoc}
\begin{ghinhoanh}\textbf{Short English to remember:}

Pride can empty applause.
\end{ghinhoanh}
\begin{tuhocnhanh}
1. Em có đang cần được công nhận quá mức không? \textit{(Are you craving recognition too much?)}

2. Thành quả nào nên để tự nó nói thay em? \textit{(Which achievement should be allowed to speak for itself?)}
\end{tuhocnhanh}
\ngancach

\section{Một người được chức lớn nhưng mất giấc ngủ.}
\begin{truyen}{Một người được chức lớn nhưng mất giấc ngủ.}{A title can cost peace.}
\chuhoa{C}{hức vụ mới đem lại tiền và uy tín, nhưng cũng kéo theo điện thoại lúc nửa đêm và áp lực triền miên. Khi nhìn mình ngồi giữa phòng khách lúc hai giờ sáng, anh hiểu có những cái được sáng rực bên ngoài mà tối hẳn bên trong.}
\textit{(The new title brought money and status, but also midnight calls and endless pressure. Sitting awake in the living room at two in the morning, he understood that some gains shine outside while darkening the inside.)}
\end{truyen}
\begin{baihoc}
Được nhiều hơn chưa chắc là sống tốt hơn, nếu phần bình yên bị lấy mất.
\textit{(Gaining more does not always mean living better, if peace is taken away.)}
\end{baihoc}
\begin{ghinhoanh}\textbf{Short English to remember:}

A title can cost peace.
\end{ghinhoanh}
\begin{tuhocnhanh}
1. Em đang theo đuổi điều gì vì vẻ ngoài của nó? \textit{(What are you chasing because of how it looks from outside?)}

2. Em có tính tới cái giá bên trong chưa? \textit{(Have you counted the inner cost yet?)}
\end{tuhocnhanh}
\ngancach

\section{Một cô gái thắng trên mạng nhưng thấy lòng trống.}
\begin{truyen}{Một cô gái thắng trên mạng nhưng thấy lòng trống.}{Winning online is not the same as living well.}
\chuhoa{C}{ô đáp trả rất hay, được nhiều người bênh và chia sẻ. Nhưng khi tắt màn hình, cô vẫn mệt, vẫn buồn, vẫn bị ám bởi cuộc hơn thua đó. Có những chiến thắng chỉ làm đông thêm người xem chứ không làm nhẹ lòng người thắng.}
\textit{(She replied brilliantly, was supported and widely shared. But after closing the screen, she still felt tired and haunted. Some victories only increase the audience, not the peace of the winner.)}
\end{truyen}
\begin{baihoc}
Không phải cái thắng nào đông người chứng kiến cũng là cái thắng đáng giữ.
\textit{(Not every victory with a large audience is worth keeping.)}
\end{baihoc}
\begin{ghinhoanh}\textbf{Short English to remember:}

Winning online is not living well.
\end{ghinhoanh}
\begin{tuhocnhanh}
1. Em có đang đem bình yên của mình đổi lấy một màn thắng nhỏ không? \textit{(Are you trading your peace for a small public win?)}

2. Việc gì nên dừng ở mạng và không nên mang vào lòng? \textit{(What should stay on the screen and not be carried into your heart?)}
\end{tuhocnhanh}
\ngancach

\section{Một người được khen là giỏi nhưng không còn ai dám góp ý.}
\begin{truyen}{Một người được khen là giỏi nhưng không còn ai dám góp ý.}{Praise can isolate.}
\chuhoa{N}{gười ta nể anh, nhưng cũng ngại nói thật với anh. Mọi việc tưởng đang tốt cho tới khi anh mắc sai lầm lớn mà không ai báo trước. Anh mới hiểu được tôn lên quá cao đôi khi cũng là bắt đầu bị cô lập.}
\textit{(People respected him, but also hesitated to tell him the truth. Everything seemed fine until a major mistake came with no warning. He then understood that being placed too high can also become a kind of isolation.)}
\end{truyen}
\begin{baihoc}
Thắng trong uy tín mà thua trong khả năng nghe sự thật thì sớm muộn cũng trả giá.
\textit{(Winning in reputation but losing the ability to hear truth will eventually cost you.)}
\end{baihoc}
\begin{ghinhoanh}\textbf{Short English to remember:}

Praise can isolate.
\end{ghinhoanh}
\begin{tuhocnhanh}
1. Ai còn đủ thật lòng để góp ý cho em? \textit{(Who still tells you the truth honestly?)}

2. Em có đang vô tình làm người khác ngại nói thật không? \textit{(Are you making others afraid to speak honestly to you?)}
\end{tuhocnhanh}
\ngancach